\section{The protocol without zero-knowledge}
\label{s:protocol:no:zk}

In this section we describe the main construction for proving a two-stage R1CS, yet without the additional measures for zero-knowledge.
The protocol is a combination Spartan \cite{Spartan} and WHIR \cite{WHIR} as a general inner-product argument. 
We skip the proving the values of the public circuit polynomials, which in \cite{Spartan} is handled by SPARK, and assume that the verifier does compute their values by itself.
In this way we obtain succinct proof sizes, yet the verifier costs are of the same order as the witness size.

%is being outsourced to the verifier, yielding only succinct proof sizes, yet not a succinct verifier.

% In the following we consider the R1CS stage 1 and 2 witnesses $\vec w_1$ and $\vec w_2$ padded to their next power of two, and hence of twoadic size $2^{m_1}$ and $2^{m_2}$.
% Without loss in generality we assume that $m_1 \geq m_2$, and we will view $w_2$ padded to the same length via lifting their MLEs.
% that $2^{m_1}$ matches $2^{k_0}$, the dimension of the largest function space $\mathscr P_{k_0}$ as specified by WHIR.


\subsection{Arithmetization.}
% The R1CS stage 1 and 2 witness vectors $\vec w_1, \vec w_2$ are zero-padded to their next power of two, and placed on the Boolean hypercube of respective size in the common manner.  
We consider the R1CS witness vectors padded to their next power of two, i.e. as functions
\[
w_i: \{0,1\}^{m_i} \longrightarrow \mathbb F_q, \quad i=1,2,
\]
%and $w: \{0,1\}^n \longrightarrow F$, 
and we denote their multilinear extensions again by the same letters.
Likewise, the number of overall constraints are padded to a power of two, say $2^m$, and we view the R1CS matrices as functions
\[
A_i, B_i, C_i: \{0,1\}^m \times \{0,1\}^{m_i} \longrightarrow \mathbb F_q, \quad i =1,2, 
\]
with the same notation for their multilinear extensions. 

Throughout the following we assume that $m_1\geq m_2$. 
% \subsubsection{Lifted R1CS.}
We lift $w_2$ to the same number of variables as $w_1$, and likewise we do for the second stage matrix polynomials $A_2, B_2, C_2$.
The resulting system has
\begin{multline}
\label{e:combined:witness:poly}
w(X_0,X_1,\ldots, X_{m_1}) =
\\
(1- X_0) \cdot w_1(X_1, \ldots, X_{m_1}) + X_0 \cdot w_2^*(X_{1},\ldots, X_{m_1}),
\end{multline}
as the combined witness polynomial, 
and
\begin{multline}
\label{e:combined:matrix:poly}
M(\vec Y,X_0,X_1,\ldots, X_{m_1}) =
\\
(1- X_0) \cdot M_1(\vec Y, X_1, \ldots, X_{m_1}) +
 X_0 \cdot  M_2^*(\vec Y, X_1,\ldots, X_{m_1}),
\end{multline}
for $M=A,B,C$, as the combined matrix polynomials.
%where $\vec Y=(Y_1,\ldots, Y_m)$, and $M_2^*(\vec Y,\:.\:)$ is the lifting of $M_2(\vec Y,\:.\:)$ to $m_1$ variables.


\ulrich{placing of inputs}
% Again, the multilinear extensions of the complete matrices $M$ and their components $M_1, M_2$ are connected via
% \begin{multline*}
% M(\vec Y, X_1,\ldots,X_n) = (1 - X_n)\cdot M_1(\vec Y, X_1, \ldots, X_{n-1}) 
% \\ + X_n \cdot M_2(\vec Y, X_1, \ldots, X_{n-1}).
% \end{multline*}



\subsection{Encoding.}

We encode $w_1$ and $w_2^*$ in interleaved manner, using the same basecode
\begin{equation}
\label{e:base:code}
C = \RS[F,D, 2^{k}],
\end{equation}
which corresponds to the first basecode of WHIR.
(We will discuss the setup of WHIR later.)
The multilinears $w_1$ and $w_2^*$ are decomposed into $(m_1 - k)$-prefix components, 
%in the $k$ variables $X_{m_1-k+1}, \ldots$, $X_{m_1}$, 
% \[
% w_{i,0}(X_{m_1-k+1},\ldots, X_{m_1}),\ldots, w_{i,e_i-1}(X_{m_1-k+1},\ldots, X_{m_1})
% \]
%as in \eqref{e:prefix:decomp}, 
which are then encoded non-systematically, using their univariate representation from $\mathscr P_{k}$.
%in the previously described values-to-coefficient regime.
% \footnotemark
% \footnotetext{%
% Given the values $(w_1(i))_{i\in H_{m_1}}$, we reshape it into a matrix of $2^{m_1-k_1}$ rows of length $2^{k_1}$, $(w_1(i + 2^{m_1-k_1}j))_{i,j}$,
% and encode each of the rows separately, taking the row entries as coefficients of the polynomials.
% }
The resulting words 
\[
f_1, f_2^* \in C^{e},
\]
the interleaved Reed--Solomon code of interleaving depth $e = 2^{m_1-k_1}$, 
where $f_2^*$ has only $d = 2^{\max\{m_2-k_1,0\}}$ non-zero components. 
(These are the prefix embeddings of the $d$ components of $w_2$.)
Thus committing to the lifted $w_2^*$ causes no overhead compared to $w_2$, given that $m_2 \geq k$.
% Whenever $m_2$ is smaller than $k$, we will work with lifted commitments.


% assuming that $m_i\geq k$. 


% This amounts to decomposing the combined witness polynomial into 
% \begin{align*}
% w(X_1,\ldots, X_{m_1 + 1}) = w(X_1,\ldots, X_k, \vec y)    
% \end{align*}



% While the setup of WHIR will be so that $m_1 \geq k$, which is typically the larger witness, it may happen that the second stage witness is smaller than the message length of $C$. 
% In this case we use \textit{lifting}, which we decribe next.

% \ulrich{%
% More elegant: Let us lift $p_2(X)$ to $p_2(X^{e})\in\mathscr P_{m_1}$ with $e=2^{m_1 - m_2}$, and decompose it into $\mathscr P_{k_1}$-components.
% If $m_2 < k_1$, i.e. $e > e_1$, we should end up with the lifting as described below:
% There $p_2^*(X) = p_2(X^{e_1 \cdot 2^{k_2 - m_2}}) $
% \[
% p_2(X^e) = p_2((X^{e_1})^{e/e_1})
% \]
% gives
% \[
% p_{2,1}(Y) = p_2(Y^{e/e_1}) = p_2(Y^{2^{m_1-m_2-m1+k_1}}) = p_2(Y^{2^{k_1 - m_2}}).
% \]
% But we do not want to commit that over $D_1$.
% }


\subsubsection{Lifting small commitments.}
Whenever $m_2 < k$, we commit the univariate representation $u_{2}(X)$ of $w_2$ over a projection of $D$, the image of the multiplicative homomorphism
\begin{equation}
\label{e:projection}
\pi: D \longrightarrow \pi(D),\quad x\mapsto x^{2^{k_1 - m_2}}.
\end{equation}
Since $D$ is an FFT domain, the map is $2^{k - m_2}$-to-$1$ and onto.
The committed word $f_2 \in \mathbb F^{D'}$ is from the smaller code
\begin{equation}
\label{e:projected:code}
C' := \RS[F, D', 2^{m_2}],
\end{equation}
which has the same relative rate as $C$.
The lifted word
\begin{equation}
\label{e:lifting:uni}
f_2^*(x) =  f_2(\pi(x)) %= f_2 (x^{2^{k-m_2}})
\end{equation}
belongs to $C$, since $u_2^*(X) = u_2(X^{2^{k-m_2}})$ is of degree less than $2^{k}$.
%
By means of the projection $\pi$ we simulate $f_2^*\in C$ without paying the price of committing over the larger domain.


\begin{rem}
As noted in Remark~\ref{rem:postfix:decomp}, the current implementation employs a postfix decomposition. However, liftings of the corresponding univariate polynomials are naturally expressed in terms of prefix embeddings, resulting in a somewhat inelegant mixture of prefix and postfix terminology.
\end{rem}

% \ulrich{Ref to Ligerito and zook}
%(See Section \ref{s:WHIR} for details.)

% We use the same commit principle on application layer (i.e. on top of WHIR) for polynomials from \textit{any} $\mathscr P = \mathscr P_k$, 
% \[
% \mathscr P_0 \supseteq \mathscr P \supsetneq \mathscr P_{k_r},
% \]
% not necessarily matching one of the spaces in \eqref{e:chain:function:spaces}: 
% %as interleaved codeword with one of the base codes $C_1,\ldots, C_r$. 
% %from one of the spaces in \eqref{e:chain:function:spaces}.
% Take the unique $i$ so that $\mathscr P_{k_i}\supseteq \mathscr P \supsetneq \mathscr P_{k_{i+1}}$,
% and denote the co-dimension of $\mathscr P_{k_{i+1}}$ in $\mathscr P$ by
% \begin{equation}
% e = \dim(\mathscr P/\mathscr P_{k_{i+1}}).
% %= 2^{k_{i} - k_{i+1}},
% \end{equation}
% The FFT-decomposition of any $p(X) \in \mathscr P$ into 
% %a vector of $e_i$ polynomials $p_1(X),\ldots, p_{e_i}(X)$ from $\mathscr P_{i+1}$,
% \begin{align}
% \label{e:FFT:decomp}
% p(X) = p_1(X^e) + X\cdot p_2(X^e) + \ldots + X^{e - 1}\cdot p_{e}(X^e),
% \end{align}
% defines a linear identity
% \begin{equation}
% \label{e:FFT:decomp:identity}
% \mathscr P \simeq \mathscr P_{k_{i+1}}^{e},    
% \end{equation}
% %
% and with this identity we commit $p(X)\in\mathscr P$ as word from 
% \begin{align}
% %\label{e:C_i}
% C_{k_{i+1}}^e 
% &= 
% \left\{
% [c_1(x),\ldots, c_{e}(x)]\big|_{x\in D_{i+1}} \::\: \text{all }c_j \in C_{k_{i+1}}
% \right\}.
% \end{align}
%\ulrich{Maybe}
% The codes $C_1,\ldots, C_r$ generated by the spaces $\mathscr P_{k_1}, \ldots \mathscr P_{k_r}$ over the domains $D_1, \ldots, D_r$, are the base codes of our scheme. 
% Their minimum distance is
% \begin{equation}
% \label{e:min:dist:baseCode}
% \delta(C_i) = 1 - \rho(C_i) - \frac{1}{|D_i|},  
% \end{equation}
% where
% \begin{equation}
% \label{e:rate:baseCode}
% \rho(C_i) = \frac{2^{k_{i}}}{|D_i|} = 2^{k_{i}- n_{i}}
% \end{equation}
% is the relative rate.
% \ulrich{%
% Maybe add blow-up factors?
% }





\subsection{The IOPP}


%Both witnesses will be committed as described above, as interleaved codeword with base code one of the $C_1,\ldots, C_r$ from \eqref{e:base:codes}.

% In many of our applications, the second stage witness is not much smaller than the first, $m_1 = m_2 + 1$, but we will not make use of this fact.


%The complete witness $\vec w = \vec w_1 \| \vec w_2$ is then of size $N=2^n$.


% The WHIR function spaces $\mathscr P_{k_0}$, \ldots $\mathscr P_{k_r}$ can be chosen freely, except that $k_0 = m_1$, so that $\mathscr P_{k_0}$ matches the size of the witness $w_1$.
% The size of $w_2$ either matches the dimension of one of these spaces, or an intermediate power of them. 
% Take $j_2 \in \{k_1, \ldots, k_r\} $ as the unique value so that $m_2 \in (k_{j_2}, k_{j_2 - 1}]$. 
% Then the univariate representation
% \[
% w_2(X) \in \mathscr P_{k_{j_2}}^{e_{j_2}'},
% \]
% with intermediate exponent $e_{j_2}' = 2^{m_2 - k_{j_2}}$ from $(1, e_{j_2}]$. 
% % 2^m_2 = e_{j_2}' \cdot 2^{k_{j_2}}
% We commit it as word from the interleaved code $C_{k_{j_2}}^{e_{j_2}'}$. 
% \ulrich{This needs to be moved to function spaces and codes.}





% \subsubsection{Commitments.}
% First stage witness $w_1$ is committed as codeword from $C_1^{e_0}$, i.e. $k_1\cdot e_1 = 2^{n_1}$ and 
% $w_2$ is committed as codeword from one of the $C_i$, with interleaving depth as in $C_i$, or less (as long as a power of two).



% The witnesses $w_1$ and $w_2$ are committed as Reed-Solomon codewords from
% \[
% C_{i} = \RS[F, D_i, 2^{n_i}] = \{p(x)|_{x\in D_i} \::\: p \in F[X]^{<2^{n_i}} \}, \quad i=1,2,
% \]
% respectively, of possible different rate, and so that their evaluation domains are WHIR compatible, i.e. $D_2 \cap \pi^{n_1-n_2}(D_1) =\varnothing$ whenever $n_2 < n_1$.
% \ulrich{%
% Not sure if this is the best setting. 
% This depends very much how we take the different withness lengths into account in WHIR. 
% I think that a constant rate, and projection-consistent domain, choice is the better way to present it here. 
% }




% \subsubsection{The proof of proximity.}
Given the commitment of $w_1$ as interleaved word $f_1$ on $D$, and proximity parameter $\theta_1\in (0,1)$, the target of the protocol is to show that $f_1$ is $\theta_1$-close to a solution of our augmented constraint system, characterized by  $\mathscr R_{main}$ and $\mathscr R_{aux}$ and the inputs $\vec{in} = (in_1,\ldots, in_{t_1})$.
%
% Recall that we use the values-to-coefficients regime, hence 
That is, we want to prove that the committed function belongs to
\begin{equation}
\label{e:IOPP:main:R:0}
\mathcal R_\theta = \left\{
    f\in \left(\mathbb F_q^{D}\right)^{e_1} \::\:
    \begin{matrix}
    \exists p(X) \in\mathscr P_{m_1},  %\in \mathbb F_q[X]^{<2^{m_1}}, 
    \: 
    d(p, f) \leq \theta_1,
    \\
    \text{ and }
    \\
    \text{the coefficients $(c_i)$ of $p(X)$ satisfy} 
    \\
    (c_i)  \in \mathscr R_{main} \cap \mathscr R_{aux},
    \\
    (c_1,\ldots,c_{t_1}) = \vec{in}
    \end{matrix}
\right\},
\end{equation}
where $d(p,f)$ refers to the interleaved representation of $p(X)$ as  $e_1 = 2^{m_1-k}$ component polynomials of degree $< 2^{k}$.
The protocol is divided into three phases.

\subsubsection{Phase 1.}
The first phase is devoted to the second stage of the R1CS.
The prover receives the stage-2 challenge $\vec\xi\sample\mathbb F_q^{t_2}$ from the verifier, and responds with the commitment $f_2$ of the witness $w_2(X_1,\ldots, X_{m_2})$ as word from the smaller code $C'$.
Via lifting \eqref{e:lifting:uni}, the committed function defines $w_2^*(X_1,\ldots, X_{m_1})$. 


\subsubsection{Phase 2.}
The Spartan zero-check.
We prove constraint satisfaction of the multilinears
\begin{align}
\label{e:a:poly}
w_M(\vec Y) &= 
\sum_{\vec x\in H_{m_1+1}} M(\vec Y, \vec x)\cdot w(\vec x), 
\\\notag
&= \sum_{\vec x\in H_{m_1}} M_1(\vec Y, \vec x)\cdot w_1(\vec x) + \sum_{\vec x\in H_{m_1}} M_2^*(\vec Y,\vec x) \cdot w_2^*(\vec x)
\\\notag
&= \sum_{\vec x\in H_{m_2}} M_1(\vec Y, \vec x)\cdot w_1(\vec x) + \sum_{\vec x\in H_{m_2}} M_2(\vec Y,\vec x) \cdot w_2(\vec x),
% \quad 
% b(\vec Y) = 
% \sum_{\vec x\in H_{m_1+ 1}} B(\vec Y, \vec x)\cdot w(\vec x), 
% \\
% c(\vec Y) =
% \sum_{\vec x\in H_{m_1+1}} C(\vec Y, \vec x)\cdot w(\vec x),
\end{align}
for $M=A,B,C$.
% \begin{equation}
% \label{e:a:poly}
% \begin{aligned}
% a(\vec Y) = 
% \sum_{\vec x\in H_{m_1+1}} A(\vec Y, \vec x)\cdot w(\vec x), 
% \quad 
% b(\vec Y) = 
% \sum_{\vec x\in H_{m_1+ 1}} B(\vec Y, \vec x)\cdot w(\vec x), 
% \\
% c(\vec Y) =
% \sum_{\vec x\in H_{m_1+1}} C(\vec Y, \vec x)\cdot w(\vec x),
% \end{aligned}
% \end{equation}
%which are uniquely determined by the commitments of $w_1$ and $w_2^*$.
%
For showing that
\begin{equation*}
% \label{e:R1CS:domain:identity}
w_A(\vec y)\cdot w_B(\vec y) - w_C(\vec y) = 0, \quad \text{for all }\vec y \in H_m.
\end{equation*}
we run the sumcheck protocol on its inner product with $eq(\:.\:,\vec\rho)$, where $\vec\rho\sample F^m$ is random.
% %of the quadratic expression $a(\vec y)\cdot b(\vec y) - c(\vec y)$ with 
% \begin{equation*}
% % \label{e:R1CS:zerocheck}
% \sum_{\vec y\in H_m} \left(a(\vec y)\cdot b(\vec y) - c(\vec y)\right) \cdot eq(\vec y, \vec \rho) = 0,
% \end{equation*}
% where $eq(\vec y,\vec\rho)$ is 
% the $m$-dimensional Lagrangian $eq(\vec y, \vec\rho)$ at a random $\vec \rho\sample F^{m}$ sampled by the verifier.
% The sumcheck protocol 
This reduces the above claim to evaluation claims 
\begin{equation}
\label{e:image:polys:evals}
v_a = w_A(\vec \alpha),\quad v_b = w_B(\vec\alpha),\quad v_c = w_C(\vec\alpha),
\end{equation}
where  $\vec\alpha\in F^m$ are the verifier challenges drawn in the course of the sumcheck rounds.
% (The value $eq(\vec\alpha, \vec \rho)$ is public and can be computed cheaply.)
%and efficiently computable in $O(m)$ steps.) 

% To prove the claims
% \begin{align*}
% v_a &= \sum_{\vec x\in H_n} A(\vec \alpha, \vec x)\cdot w(\vec x), \quad v_b = \sum_{\vec x\in H_n} B(\vec \alpha, \vec x)\cdot w(\vec x)
% \\
% v_c &= \sum_{\vec x\in H_n} C(\vec \alpha, \vec x)\cdot w(\vec x),
% \end{align*}
% together with the consistency constraints with respect to the public circuit inputs, including constants and the logUp challenges,
% \begin{align*}
% \vec {in}_j &= \sum_{\vec x\in H_n} eq(j, \vec x)\cdot w(\vec x), \quad j=0,\ldots,I,
% \\
% \alpha_j &= \sum_{\vec x\in H_n} eq(2^{n-1} + j,\vec x) \cdot w(\vec x), \quad j= 0,\ldots,k
% \end{align*}

\subsubsection{Phase 3.}
The third phase of the protocol is devoted to proving the linear claims \eqref{e:image:polys:evals}, which is done by WHIR on the combined witness polynomial \eqref{e:combined:witness:poly}.
For public input and challenge consistency, we further add the constraints
\begin{align}
\label{e:input:w1:constraint}
in_j &= \sum_{\vec x\in H_{m_1}} eq(j,\vec x)\cdot w_1(\vec x), \quad j=0,\ldots, t_1,
\\
\label{e:input:w2:constraint}
\xi_j & = \sum_{\vec x\in H_{m_1}} eq^*(j,\vec x)\cdot w_2^*(\vec x), \quad j=0,\ldots, t_2,
\end{align}
where $\vec\xi = (\xi_1,\ldots, \xi_{t_2})$ are the challenges drawn in the second stage of the R1CS.
See Section \ref{s:WHIR} for the description of WHIR.
Assuming that the protocol is configured to consume every variable of $w$, 
the phase eventually reduces the linear claims on $w$ to claims on the combined matrix polynomials, 
\begin{equation}
\label{e:split:matrix:values}
A(\vec\alpha,\vec\beta), 
\quad 
B(\vec\alpha,\vec\beta), 
\quad 
C(\vec\alpha,\vec\beta),
\end{equation}
where $\vec\beta \in F^{m_1}$ are the folding randomnesses drawn in the course of WHIR. 

Note that computing the matrix values \eqref{e:split:matrix:values} requires $O(n)$ field operations, where $n$ is the maximum number of non-zero entries in the R1CS matrices. 
While these values can be proven via the SPARK polynomial commitment scheme for sparse multilinears \cite{Spartan}, our main use case pivots for letting the verifier take the burden of matrix evaluation

\subsubsection{Summary.}
Let us summarize the steps of our IOPP. 
For simplicity we assume that 
\[
t_1 = t_2 = t.
\]
Recall that $m_1\geq m_2$, which is only for simplicity of our presentation. 
The protocol as well as its security properties does not change if otherwise.

% \begin{equation}
% \label{e:split:matrix:values}
% A(\vec\alpha,\vec\beta), 
% \quad 
% B(\vec\alpha,\vec\beta), 
% \quad 
% C(\vec\alpha,\vec\beta),
% \end{equation}
%and those of the involved Lagrangians $L_z(\vec\beta)$ and $eq(j,\vec \beta)$, 


% As we emphasized before, we do not target a succinct verifier and leave the verification of \eqref{e:split:matrix:values} to the verifier, instead of providing a separate proof for it.


\begin{protocol}[IOPP without zk]
\label{prot:main}
The R1CS matrices $A_i,B_i,C_i \in \mathbb F_q^{m\times m_i}$, $i=1,2$, and public inputs $\vec{in} = (in_0,\ldots, in_{t})\in \mathbb F_q^t$ are known to both the prover and the verifier. 
The prover has committed to the first-stage witness $\vec w_1\in \mathbb F_q^{m_1}$ as interleaved Reed--Solomon codeword $f_1\in C^e$ as described above, and provided the verifier oracle access.
\begin{enumerate}
\item
\label{i:main:phase:1}
R1CS second stage:
% \begin{enumerate}
% \item
% \label{i:main:DEEP:w1}
% The verifier sends a random $z_1\sample F$ to the prover, which answers with the value $v_1$ of the univariate representation $w_1(X)$ at $z_1$.
% %
% \item
% \label{i:main:logUp}
The verifier samples the stage-2 challenge
\[
\vec\xi \sample \mathbb F_q^{t}
\]
and sends them to the prover.
The prover computes $\vec w_2 \in \mathbb F_q^{m_2}$  the second-stage witness of the R1CS, commits to it as codeword $f_2\in C'$, as described above, and provides the verifier oracle access.
%(the logUp inverses and their sums), 
% commits to it as interleaved Reed--Solomon codeword $f_2$ associated with $\mathscr P_{m_2}(\mathbb F_q)$, and provides the verifier oracle access. 
%
% \item 
% \label{i:main:DEEP:w2}
% The verifier sends a random $z_2\sample F$ to the prover, which answers with the value of the univariate representation $w_2(X)$ at $z_2$.
% \end{enumerate}
%answers with the logUp R1CS witnesses $\vec w_2$, so that the 

\item 
\label{i:main:phase:2}
R1CS zero-check:
% Both prover and verifier engage in a \textit{zero-check} for the domain identity
% \begin{equation}
% a(\vec x) \cdot b(\vec x) = c(\vec x), \quad \vec x\in H_n = \{0,1\}^n
% \end{equation}
% where $a, b$ and $c$ are the multilinear extensions of the implicitly committed matrix-images
% \[
% \vec a = A\cdot \vec w, \vec b = B\cdot \vec w, \vec c = C\cdot \vec w.
% \]
The verifier sends a random 
\[
\vec\rho \sample F^m,
\]
and both prover and verifier engage in the sumcheck protocol for %the R1CS constraints
\begin{align}
\label{e:R1CS:zerocheck}
\sum_{\vec x\in H_m} \left(w_A(\vec x)\cdot w_B(\vec x) - w_C(\vec x)\right)\cdot eq(\vec x, \vec\rho) = 0,
\end{align}
where $w_A(\vec x), w_B(\vec x), w_C(\vec x)$ are the multilinears as specified in \eqref{e:a:poly}.
After the last round of the sumcheck, the prover sends 
% (The protocol reduces the claim on the hypercube sum \eqref{e:R1CS:zerocheck} to a quadratic check on the evaluation claims 
\begin{equation*}
v_A = w_A(\vec\alpha),  \quad v_B = w_B(\vec\alpha), \quad v_C = w_C(\vec\alpha),
\end{equation*}
where $\vec\alpha \in F^m$ are the verifier challenges drawn in the course of the sumcheck, and the verifier checks if 
$q_m(\alpha_m) = (v_A\cdot v_B - v_C) \cdot eq(\alpha,\rho)$.
%These claims are announced by the prover in the last round. 

\item 
\label{i:main:phase:3}
WHIR: 
% The verifier samples a random $\lambda\sample F$, 
Both prover and verifier run WHIR (Section \ref{s:WHIR}) for the following linear constraints on the combined witness polynomial $w$:
\begin{align}
\label{e:phase3:claim:vM}
v_M = \sum_{\vec x\in H_{m_1}} M_1(\vec\alpha, \vec x) \cdot w_1(\vec x) + \sum_{\vec x\in H_{m_1}} M_2^*(\vec\alpha,\vec x)\cdot w_2^*(\vec x)
\end{align}
with $M = A,B,C$, and
\begin{align}
\label{e:phase3:claim:inputs}
in_i = \sum_{\vec x \in H_{m_1}} eq(i, \vec x) \cdot w_1(\vec x),
\quad 
\xi_i = \sum_{\vec x \in H_{m_2}} eq^*(i, \vec x) \cdot w_2^*(\vec x),
\end{align}
for $i = 0,\ldots, t$.
% (The concrete form of the linear combination is described below.)
% (This phase reduces the validity of the linear constraints to claims on the combined matrix polynomials
% \begin{equation}
% \label{e:split:matrix:values}
% A(\vec\alpha,\vec\beta), 
% \quad 
% B(\vec\alpha,\vec\beta), 
% \quad 
% C(\vec\alpha,\vec\beta),
% \end{equation}
% where $\vec\beta \in F^{m_1}$ are the folding randomnesses drawn in the course of WHIR.)
\end{enumerate}
If all verifier checks in Phase \ref{i:main:phase:2} and Phase \ref{i:main:phase:3}, including that of the final matrix values \eqref{e:split:matrix:values}, pass, then the verifier accepts.
Otherwise, it rejects.
\end{protocol}

% The three claims \eqref{e:phase3:claim:vM} and \eqref{e:phase3:claim:inputs} are of inner products of $w$ 
% with the kernels
% \[
% K_M = (1-X_0) \cdot M_1(\vec\alpha,X_1,\ldots, X_{m_1}) + X_0\cdot M_2^*(X_1,\ldots, X_{m_1}), 
% \]
% $M=A,B,C$, and
% \begin{align*}
% K_{1,i} &= (1-X_0)\cdot  eq(i,X_1,\ldots, X_{m_1}), 
% \\
% K_{2,i} &= X_0\cdot eq^*(i,X_1,\ldots, X_{m_1}),
% \end{align*}
% for $i=0,\ldots, t$.
%The three claims \eqref{e:phase3:claim:vM} and \eqref{e:phase3:claim:inputs}  are combined into a single linear claim on $w$, 

\begin{rem}
\label{rem:WHIR:not all variables consumed}
If WHIR is configured to consume only $m_1-k_r$ variables, for some small $k_r$, then the matrix polynomials in the final WHIR claim are specialized to $M(\vec\alpha, \vec\beta\|\:.\:)$, $M=A,B,C$, which are multilinears in the remaining $k_r$ variables. 
\end{rem}



\subsection{Soundness.}
\label{s:IOPP:soundness:discussion}

Without specifying the formal security notions, we state the soundness guarantees of the protocol, for proximity parameters up to the Johnson radius.
This means that $\theta$ as in the main relation \eqref{e:IOPP:main:R:0} satisfies
\[
\theta_1 \in \left(0, 1 -\sqrt{1 - \delta_1}\right),
\]
where $\delta_1$ is the relative minimum distance of the base code $C_1 = C$. 
The corresponding list size bound is
\[
\ell_1(\theta) =  \frac{1}{2 \sqrt{1-\delta_1}} \cdot \frac{1}{\eta},
\]
the combinatorial Johnson bound, with  $\eta =1 - \sqrt{1 - \delta_1} - \theta$.
WHIR with $r$ rounds involves the further Reed--Solomon codes $C_2,\ldots, C_r$, see Section \ref{s:WHIR} below, with each round having its own proximity parameter $\theta_i$, again assumed to be within the Johnson radius. 

\begin{thm}
\label{thm:main:soundness:informal}
For proximity parameters $\theta_i\in \left(0, 1 - \sqrt{1-\delta_i} \right)$, $i=1,2,\ldots, r$, where $\delta_i$ are the minimum distances of the Reed--Solomon codes $C_1$, \ldots, $C_r$.
Protocol \ref{prot:main} is a round-by-round knowledge sound interactive oracle proof for that $f_1\in \mathcal R_{\theta_1}$ as specified in \eqref{e:IOPP:main:R:0}, with the following round-by-round soundness error for the three phases:
\begin{align}
\label{e:phase1:soundness:error}
\varepsilon_\text{phase 1} 
&\leq 
    % % \max\left( 
    % \binom{\ell_ 1}{2} \cdot \frac{2^{m_1}}{|F|} %\quad
    % + \ell_2\cdot \varepsilon_{stage 2} %\quad
    % + \binom{\ell_2}{2} \cdot \frac{2^{m_2}}{|F|},
    % % \max\right)
    \ell_1(\theta_1) \cdot \varepsilon_\text{R1CS} 
\\\intertext{%
where $\varepsilon_{R1CS}$ is the soundness error of the randomized AIR, 
}\label{e:phase2:soundness:error}
\varepsilon_\text{phase 2}  
&\leq \ell_1(\theta_1) \cdot \max\left\{\frac{m}{|F|}, \frac{3}{|F|}\right\},
\end{align}
and 
\begin{equation}
\varepsilon_\text{phase 3} \leq 
%\ell_1\cdot \frac{3 + 2t - 1}{|F|} + 
\varepsilon_\text{WHIR}(\theta_1,\ldots, \theta_r)
\end{equation}
for the round-by-round soundness error of WHIR with $r$ rounds on the inner product claims \eqref{e:phase3:claim:vM} and \eqref{e:phase3:claim:inputs}.
\end{thm}



% \begin{rem}
% Whenever the commitments $f_1$, $f_2$ belong to the same base code, and hence are absorbed in the same folding step of WHIR, the quadratic behavior in the list size is reduced to $\varepsilon_\text{phase 1} \leq \ell_1 \cdot \varepsilon_{stage 2}$ and $\varepsilon_\text{phase 2} \leq \ell_1\cdot \left(\frac{m}{|F|} + m\cdot \frac{3}{|F|}\right)$, since in this case, WHIR proves correlated agreement.
% See the discussion below.
% \end{rem}

% The soundness analysis behind Theorem \ref{thm:soundness:main:informal} is largely standard.
% Nevertheless we quickly discuss the soundness errors of the phases, and how lifting plays into it.
% (The role of lifting)
% \begin{equation}
% \label{e:IOPP:main:R:0}
% \mathcal R_\theta = \left\{
%     f_1\in \left(\mathbb F_q^{D}\right)^{e} \::\:
%     \begin{matrix}
%     \exists p(X) = \sum_{i=0}^{2^{m_1}- 1} c_i X^i,  %\in \mathbb F_q[X]^{<2^{m_1}}, 
%     d(p, f) \leq \theta,
%     \\
%     \text{ with }
%     \\
%     (c_i)  \in \mathscr R_{main} \cap \mathscr R_{aux},
%     \\
%     (c_1,\ldots,c_{t_1}) = \vec{in}
%     \end{matrix}
% \right\},
% \end{equation}

Let us roughly sketch the analysis behind Theorem \ref{thm:main:soundness:informal}.
First of all, the entire protocol aims to prove a proximate solution of the \textit{lifted} R1CS, specified by $A_1, B_1, C_1$ and the prefix-embedded second stage matrices $A_2^*, B_2^*, C_2^*$.
%which have the same number of columns as the first stage matrices.
%padded to the same number of columns as the first stage matrices.
% Unlike in \cite{lifted:FRI},
Unlike in certain univariate arithmetizations, see \cite{liftedFRI}, lifting by prefix-embedding does not change the circuit logic -- it only introduces variables that are never used.
The lifted system has the same soundness error as the original one, i.e. $\varepsilon_{R1CS^*} = \varepsilon_{R1CS}$, and from the point of view of $w_1$, it is equivalent to ask for a second stage response satisfying the lifted system, or a smaller witness satisfying the original system.\footnotemark
\footnotetext{%
In fact, a solution of the non-lifted system is simply obtained by truncation.
}
This is the reason why Phase \ref{i:main:phase:1} is a randomized reduction from the main claim $f_1\in \mathcal R_{\theta_1}$ to that the extended prover message $(f_1, f_2^*)$ belongs to 
\begin{align}
\label{e:}
\mathcal R_\text{phase 1} 
&=
\left\{
    \begin{matrix}
        f_1 \in \left(\mathbb F_q^{D}\right)^{e_1},% v_1\in F 
        \\
        f_2^* \in \left(\mathbb F_q^{D}\right)^{e_1}%, v_2\in F 
    \end{matrix}
    \::\:
    % \hspace*{-0.25cm}
    \begin{matrix}
        \exists (p_1(X),p_2(X)) \in \mathscr P_{m_1}^2 %\times \mathscr P_{m_1}
        \\
        d((p_1,p_2), (f_1,f_2^*)) \leq \theta_1,
        % \\
        % p_1(z_1) = v_1, p_2(z_2) = v_2.
        \\
        \text{and}
        \\
        \text{their coefficient vectors}
        \\
        (\vec c_1, \vec c_2) \text{ satisfy the lifted R1CS }
        \\
        \text{on inputs }\vec{in} , \vec\xi 
    \end{matrix}
\right\}.
\end{align}
(Here, the distance refers to the interleaved representation of $p_1$ and $p_2$.)
%The latter will be ensured by Phase 2 and Phase 3 of the protocol.
Note, that the proximate $p_2(X)$ is not demanded of lifted form, in other words, it is not asked to stem from a $\theta_1$-proximate of the non-lifted word $f_2$ on $D'$.
(In fact, we do not rule out that $f_2$ is not as close to $C'$.)
A usual pigeon-hole argument on the number of $\theta_1$-proximates of $(f_1,f_2^*)$ bounds the soundness error of this reduction by
\[
\ell_1(\theta_1) \cdot \varepsilon_{R1CS^*} = \ell_1(\theta_1)\cdot \varepsilon_{R1CS},
\]
as stated in the theorem.
% Having a solution to the lifted R1CS is 

% It simply demands any proximate $(p_1(X), p_2(X))$ of $(f_1,f_2^*)$ which solves the \textit{lifted} R1CS,  however, truncation  automatically yields a solution of the non-lifted constraint system. 
% Thus we may use $\varepsilon_{R1CS}$ of the non-lifted system to argue the soundness error of Phase 1, scaled by $\ell_1(\theta_1)$ due to the usual pigeon-hole principle.
% Therefore, if the prover is able to respond on more than $\ell_1(\theta_1)\cdot \varepsilon_\text{R1CS}\cdot |F^t|$ many challenges $\vec\xi$ with a function $f_2^*$ so that $(f_1, f_2^*)\in \mathscr R_\text{phase 1}$, then by the pigeon-hole principle, at least one of the $\theta_1$-proximates of $f_1$ must have more than $\varepsilon_\text{R1CS}\cdot |F^t|$ many continuations which solves the non-lifted R1CS.

% The error for Phase 1 is the soundness error of randomized R1CS, scaled by the number of proximates of the committed functions. 
% Since both $f_1$ and $f_2$ are committed over the same base code, we view them as elements of the extended interleaved code
% \[
% (f_1,f_2)\in C_1^{2e},
% \]
% and $\ell_1$ is the list size bound at distance $\theta_1$, inherited from the base code $C_1$.
% WHIR will enforce correlated agreement between the components of the two, formally words of the interleaved code $C_1^{2e}$. 
% In general, $f_1$ and $f_2$ belong to different base codes, and we assume all combinations of the proximates,
% \[
% \ell_1 \cdot \ell_2
% \]
% as a crude upper bound for the proximates of the pair of commitments.
% Although we believe that a cross-domain analysis in the style of \cite{Stwo:whitepaper} can be extended to WHIR, thereby reducing this term, we do not pursue this direction here due to the substantial technical challenges involved.
% Our extension field $F$ provides enough security margin to handle quadratic (or even cubic) products of list sizes.
The same list size bound also scales the soundness errors of Phase \ref{i:main:phase:2}, which is a reduction from $\mathcal R_\text{phase 1}$ to that the prover messages $f_1,f_2^*$, extended by its proposed sumcheck polynomials $q_1(X),\ldots, q_m(X)$ and the final claims $v_A, v_B, v_C$, belongs to 
\begin{equation}
\label{e:}
\mathcal R_\text{phase 2} =
\left\{
    \begin{matrix}
        f_1 \in \left(\mathbb F_q^{D}\right)^{e_1}, %v_1\in F 
        \\
        f_2^* \in \left(\mathbb F_q^{D}\right)^{e_1}, %v_2\in F 
        \\
        q_1(X), \ldots, q_m(X) \in F[X]^{\leq 3}
        \\
        v_A, v_B, v_C \in F
    \end{matrix}
    \::\:
    \hspace*{-10mm}
    \begin{matrix}
        \exists (p_1(X),p_2(X)) \in \mathscr P_{m_1}^2
        \\
        d((p_1,p_2), (f_1,f_2^*)) \leq \theta_1,
        \\\text{and}
        \\
        0 = q_1(0) + q_1(1),
        \\
        q_i(\alpha_i) = q_{i+1}(0) + q_{i+1}(1),
        \\
        \text{for }i = 1,\ldots, m-1,
        \\
        q_m(\alpha_m) = (v_A \cdot v_B - v_C)\cdot eq(\vec\alpha, \vec\rho)
    \end{matrix}
\right\}.
\end{equation}
The first term of the soundness error \eqref{e:phase2:soundness:error},
\[
\ell_1(\theta_1) \cdot \frac{m}{|F|}
\]
attributes the reduction from the domain identity to the inner product with $eq(\:.\:,\vec\rho)$.
%and corresponds to the probability that one of the proximates yields a non-zero constraint multilinear which evaluates to zero at the given point $\vec\rho$.
The other term, 
\[
\ell_1(\theta_1)\cdot \frac{3}{|F|},
\]
is the soundness errors of each of the $m$ rounds of the sumcheck protocol, on the cubic zero-check expression.

Finally, and without going into details, Phase \ref{i:main:phase:3} is a reduction from $\mathcal R_\text{phase 2}$ to the final verifier relation $\mathcal R_{V}$, which trades the correlated agreement claim
\[
\exists (p_1(X),p_2(X)) \in \mathscr P_{m_1}^2, \quad
        d((p_1,p_2), (f_1,f_2^*)) \leq \theta_1,
\]
for consistency checks on the WHIR sumcheck polynomials and the aggregated WHIR claim on the final polynomial revealed in plain. 
% (The latter claim involves the evaluation claims on the matrix polynomials $M(\vec\alpha, \vec \beta)$.)
\begin{align}
\mathcal R_{V} =
\left\{
    \begin{matrix}
        f_1 \in \left(\mathbb F_q^{D}\right)^{e}, %v_1\in F 
        \\
        f_2^* \in \left(\mathbb F_q^{D}\right)^{e}, %v_2\in F 
        \\
        q_1(X), \ldots, q_m(X) \in F[X]^{\leq 3},
        \\
        v_A, v_B, v_C \in F,
        \\
        q_{1,1}(X),\ldots, q_{1,\epsilon_1}(X) \in F[X]^{\leq 3},
        \\
        g_2 \in \left(\mathbb F_q^{D_{2}}\right)^{e_{2}},  u_2\in F, 
        \\
        \vdots
        % \\
        % q_{r-1,1}(X),\ldots, q_{r-1,\epsilon_1}(X) \in F[X]^{\leq 3},
        % \\
        % g_{r-1} \in \left(\mathbb F_q^{D_{r}}\right)^{e_{r}}, u_{r-1}\in F,
        % \\
        \\
        q_{r,1}(X),\ldots, q_{r,\epsilon_r}(X) \in F[X]^{\leq 3},
        \\
        g_r \in \mathscr P_{k_r} 
    \end{matrix}
    \::\:
    % \hspace*{-5mm}
    \begin{matrix}
        % \exists (p_1(X),p_2(X)) \in \mathscr L_{m_1}\big|_{z_1,v_1}\times \mathscr L_{m_2}\big|_{z_1,v_1}
        % \\
        % d(p_1, f_1) \leq \theta_1, d(p_2, f_2)\leq \theta_2
        % % \\
        % % p_1(z_1) = v_1, p_2(z_2) = v_2,
        % \\\text{and}
        \\
        0 = q_1(0) + q_1(1),
        \\
        q_i(\alpha_i) = q_{i+1}(0) + q_{i+1}(1),
        \\
        \text{for }i = 1,\ldots, m-1,
        \\
        q_m(\alpha_m) = (v_A \cdot v_B - v_C)\cdot eq(\vec\alpha, \vec\rho)
        \\\text{and}
        \\
        \text{the WHIR consistency checks hold}
        \\
        \text{and}
        \\
        g_r \text{ satisfies the aggregated claim}
        \\
        \bar v = \langle K_r, g_r\rangle_{H_{k_r}}, 
        \\
        \text{which involves the specializations}
        \\
        M_i(\vec\alpha, \vec\beta \| \: .\:), 
        \\ 
        M=A,B,C, \: i=1,2
    \end{matrix}
\right\}.
\end{align}
Its soundness error is $\varepsilon_\text{WHIR}(\theta_1,\ldots, \theta_r)$, the overall error of WHIR with $r$ rounds on the five claims \eqref{e:phase3:claim:vM} and \eqref{e:phase3:claim:inputs}.
We describe the protocol in the next section.